% GENERATED, edit the French unit or the translation table
\mxunit{automatismes-ch04}{Functions: representations and reference functions}

\begin{fluencysheet}{Thu~3~Dec.~2026}{Real-valued function; domain of definition}
\mxfq{Warm-up · Quick}{EQ}{Solve $x-4=0$ then $2x+6=0$.}{$x=4$; $x=-3$.}{}
\mxfq{Core · Method}{MOD}{Domain of definition of $f(x)=\dfrac{1}{x-4}$?}{Exclude the value that makes the denominator zero: $\R\setminus\{4\}=\intoo{-\infty}{4}\,\cup\,\intoo{4}{+\infty}$.}{}
\mxfq{Core + · MCQ}{MOD}{The domain of definition of $h(x)=\dfrac{1}{x^2+1}$ is:\mxchoices{$\R\setminus\{1\}$}{$\R\setminus\ens{-1;1}$}{$\R$}{$\R^*$}}{C.\ $\R$}{$x^2+1\geqslant1>0$: the denominator is never zero.}
\mxfq{Review · Table}{PNP}{Is this table a proportionality table?\par\smallskip\mxarraystretch{1.15}\begin{tabular}{|c|c|c|c|}\hline
$x$ & 2 & 5 & 8 \\ \hline
$y$ & 6 & 15 & 24 \\ \hline
\end{tabular}\par }{Yes: $\dfrac{y}{x}=3$ for each column.}{}
\mxfq{Thread · Code reading}{ALG}{What does $f(4)$ return?\mxpseudo{def f(x):\\ \mxtab return 3*x - 2}}{$10$.}{}
\end{fluencysheet}

\begin{fluencysheet}{Mon~7~Dec.~2026}{Curve with equation $y=f(x)$; images and preimages}
\mxfq{Warm-up · Quick}{IMG}{$f(x)=-x^2+4$. Calculate $f(0)$, $f(-2)$ and $f(3)$.}{$4$; $0$; $-5$.}{}
\mxfq{Core · Graph}{IMG}{Curve of $f$:\par\smallskip\begin{tikzpicture}[scale=0.45,>=stealth]
\draw[very thin,gray!45] (-3,-5) grid (3,4);
\draw[->] (-3,0)--(3.3,0); \draw[->] (0,-5)--(0,4.3);
\foreach \i in {-3,-2,-1,1,2,3}\draw (\i,.1)--(\i,-.1) node[below,font=\tiny]{\i};
\foreach \j in {-5,-4,-3,-2,-1,1,2,3,4}\draw (.1,\j)--(-.1,\j) node[left,font=\tiny]{\j};
\draw[thick,mxink,domain=-3:3,samples=60,smooth] plot (\x,{-\x*\x+4});

\end{tikzpicture}\par Read off $f(1)$ and the preimages of $0$.}{$f(1)=3$; preimages of $0$: $-2$ and $2$.}{Expression used for the plot: $f(x)=-x^2+4$.}
\mxfq{Core + · Thinking}{IMG}{For this function, does the equation $f(x)=5$ have any solutions?}{No: $f(x)=4-x^2\leqslant4<5$ for all $x$; the curve does not rise as far as $y=5$.}{}
\mxfq{Review · Quick}{PRB}{Urn: $5$ red balls, $3$ blue, $2$ green. Probability of drawing a blue one? A non-red one?}{$\dfrac{3}{10}$; $\dfrac{5}{10}=\dfrac12$.}{}
\mxfq{Thread · Method}{EQ}{Solve in $\R$: $4-3x=10$.}{$x=-2$}{Pitfall: from $-3x=6$, writing $x=2$ instead of $x=\dfrac{6}{-3}=-2$.}
\end{fluencysheet}

\begin{fluencysheet}{Tue~8~Dec.~2026}{Reference functions: absolute value, square, inverse}
\mxfq{Warm-up · Quick}{CAL}{Calculate $|{-}\num{2,5}|$, $(-3)^2$ and $\dfrac{1}{-4}$.}{$\num{2,5}$; $9$; $-\num{0,25}$.}{}
\mxfq{Core · Table}{REG}{Complete, for $x=-2\,;\,-1\,;\,0\,;\,1\,;\,2$, the values of $x\mapsto|x|$, $x\mapsto x^2$ and $x\mapsto\dfrac1x$.}{$|x|$: $2\,;\,1\,;\,0\,;\,1\,;\,2$; $x^2$: $4\,;\,1\,;\,0\,;\,1\,;\,4$; $\dfrac1x$: $-\dfrac12\ ;\ -1\ ;\ \text{undefined}\ ;\ 1\ ;\ \dfrac12$.}{}
\mxfq{Core + · MCQ}{REG}{Which reference curve passes through $\pt{-1;1}$, $\pt{0;0}$, $\pt{1;1}$ and $\pt{2;4}$?\mxchoices{the line $y=x$}{the inverse function}{the absolute value function}{the square function}}{D.\ the square function}{The absolute value function would pass through $\pt{2;2}$; the inverse function is not defined at $0$; the line $y=x$ passes through $\pt{-1;-1}$.}
\mxfq{Review · Quick}{PRO}{An object is enlarged with scale factor $\num{1,5}$. What does a length of $8$ cm become? By what number is the area multiplied?}{$12$ cm; the area is multiplied by $\num{1,5}^2=\num{2,25}$.}{}
\mxfq{Thread · Prove}{LIT}{Conjecture: ``the sum of two consecutive integers is odd''. Prove it.}{$n+(n+1)=2n+1$, of the form $2k+1$ with $k=n$: odd.}{}
\end{fluencysheet}

\begin{fluencysheet}{Thu~10~Dec.~2026}{Square root and cube; sign of reference functions}
\mxfq{Warm-up · Quick}{CAL}{Calculate $\sqrt{49}$, $\sqrt{\num{0,25}}$ and $(-2)^3$.}{$7$; $\num{0,5}$; $-8$.}{}
\mxfq{Core · Method}{MOD}{Study the sign of $f(x)=-2x+6$.}{$-2x+6=0\Leftrightarrow x=3$; $f(x)>0$ for $x<3$, $f(x)<0$ for $x>3$ (negative gradient).}{}
\mxfq{Core + · True/false}{MOD}{True or false? Justify. For every real number $x$, $\sqrt{x^2}=x$.}{False. Counter-example: $x=-3$: $\sqrt{9}=3\neq-3$. In fact $\sqrt{x^2}=|x|$.}{}
\mxfq{Review · Quick}{PER}{How many squares make up a net of a cube? Does a cross (a row of $4$ squares + $1$ square on each side) work?}{$6$ squares; yes, the cross folds into a cube.}{}
\mxfq{Thread · Quick}{SOL}{How many faces, vertices and edges does a square-based pyramid have?}{$5$ faces, $5$ vertices, $8$ edges.}{}
\end{fluencysheet}

\begin{fluencysheet}{Mon~14~Dec.~2026}{Graphical and numerical solving: $f(x)=k$, $f(x)<k$}
\mxfq{Warm-up · Quick}{EQ}{Solve $x^2=16$ then $x^2=-4$.}{$x=\pm4$; no real solution ($x^2\geqslant0$).}{}
\mxfq{Core · Graph}{REG}{Curve of $f(x)=x^2$:\par\smallskip\begin{tikzpicture}[scale=0.45,>=stealth]
\draw[very thin,gray!45] (-3,-1) grid (3,6);
\draw[->] (-3,0)--(3.3,0); \draw[->] (0,-1)--(0,6.3);
\foreach \i in {-3,-2,-1,1,2,3}\draw (\i,.1)--(\i,-.1) node[below,font=\tiny]{\i};
\foreach \j in {-1,1,2,3,4,5,6}\draw (.1,\j)--(-.1,\j) node[left,font=\tiny]{\j};
\draw[thick,mxink,domain=-2.45:2.45,samples=60,smooth] plot (\x,{\x*\x});

\end{tikzpicture}\par Solve graphically $f(x)=4$ then $f(x)<4$.}{$f(x)=4$: $x=-2$ or $x=2$; $f(x)<4$: $x\in\intoo{-2}{2}$.}{Look for the abscissas of the points of the curve lying below the line $y=4$.}
\mxfq{Core + · Method}{EQ}{Solve $|x|<3$ then $|x|=5$.}{$\intoo{-3}{3}$; $x=-5$ or $x=5$.}{}
\mxfq{Review · Pitfall}{SUB}{Calculate $-x^2$ then $(-x)^2$ for $x=-3$.}{$-(-3)^2=-9$ and $(-(-3))^2=9$.}{}
\mxfq{Thread · Method}{DEM}{In a triangle $ABC$, $M$ and $N$ are the midpoints of $[AB]$ and $[AC]$. What is being proved? Which property is cited?}{$(MN)\parallel(BC)$ and $MN=\dfrac12BC$: the midpoint theorem.}{}
\end{fluencysheet}

\begin{fluencysheet}{Tue~15~Dec.~2026}{Solving $f(x)=g(x)$ and $f(x)<g(x)$}
\mxfq{Warm-up · Quick}{EQ}{Solve $2x-1=x+3$.}{$x=4$.}{}
\mxfq{Core · Graph}{REG}{Curves of $f(x)=x$ and $g(x)=x^2$:\par\smallskip\begin{tikzpicture}[scale=0.45,>=stealth]
\draw[very thin,gray!45] (-1,-1) grid (2,4);
\draw[->] (-1,0)--(2.3,0); \draw[->] (0,-1)--(0,4.3);
\foreach \i in {-1,1,2}\draw (\i,.1)--(\i,-.1) node[below,font=\tiny]{\i};
\foreach \j in {-1,1,2,3,4}\draw (.1,\j)--(-.1,\j) node[left,font=\tiny]{\j};
\draw[thick,mxink,domain=-1:2,samples=60,smooth] plot (\x,{\x});
\draw[thick,mxink,domain=-1:2,samples=60,smooth] plot (\x,{\x*\x});

\end{tikzpicture}\par Solve $f(x)=g(x)$ then $f(x)<g(x)$.}{$x=x^2\Leftrightarrow x(x-1)=0$: $x=0$ or $x=1$; $f(x)<g(x)$ on $\R$: $x(x-1)>0$ (product of two factors of the same sign), so the parabola is above the line for $x<0$ or $x>1$ (on the figure, drawn on $\intff{-1}{2}$, one reads only $\intfo{-1}{0}$ and $\intof{1}{2}$).}{Intervals: $\intoo{-\infty}{0}\cup\intoo{1}{+\infty}$.}
\mxfq{Core + · Code reading}{ALG}{What does \code{f(-1) == g(-1)} return? And \code{f(0) == g(0)}?\mxpseudo{def f(x):\\ \mxtab return x*x\\ def g(x):\\ \mxtab return 2*x + 3}}{\code{True} ($1=-2+3$) then \code{False} ($0\neq3$): $-1$ is a solution of $f(x)=g(x)$, $0$ is not.}{}
\mxfq{Review · Problem}{MOD}{Plan A: $\num{15000}$ Ar fixed + $\num{1000}$ Ar per session. Plan B: $\num{2500}$ Ar per session. From how many sessions is A cheaper?}{$\num{15000}+\num{1000}x<\num{2500}x\Leftrightarrow x>10$: from $11$ sessions ($10$: equal at $\num{25000}$ Ar).}{}
\mxfq{Thread · Quick}{REP}{Calculate the coordinates of the midpoint of $[AB]$ with $A\pt{-2;5}$ and $B\pt{6;1}$.}{$\left(\dfrac{-2+6}{2};\dfrac{5+1}{2}\right)=\pt{2;3}$.}{}
\end{fluencysheet}
